
\exo

\medskip

\begin{minipage}{0.5\linewidth}
 On donne la figure ci-dessous :
 
 \medskip
 
\begin{center}
\psset{unit=0.7cm,algebraic=true,dimen=middle,dotstyle=o,dotsize=5pt 0,linewidth=1pt,arrowsize=3pt 2,arrowinset=0.25}
\begin{pspicture*}(-1,-1)(9,5)
\psgrid[gridlabels=0,subgriddiv=1,griddots=10](-1,-1)(9,5)
\psframe[linewidth=1pt](4,0)(8,4)
\psframe[linewidth=1pt](0,0)(4,4)
\psline[linewidth=1pt](0,4)(4,0)
\psline[linewidth=1pt](0,0)(4,4)
\psline[linewidth=1pt](8,4)(4,0)
\psline[linewidth=1pt](4,4)(8,0)
\rput[br](0,4.2){$F$}
\rput[b](4,4.2){$E$}
\rput[bl](8,4.2){$D$}
\rput[t](2,1.8){$G$}
\rput[t](6,1.8){$H$}
\rput[tr](0,-0.2){$A$}
\rput[t](4,-0.2){$B$}
\rput[tl](8,-0.2){$C$}
\end{pspicture*}
\end{center}
 
 \end{minipage}
\begin{minipage}{0.5\linewidth}
Compléter :

\medskip

%\begin{multicols}{2}
\begin{enumerate}%[label=\bf\alph{{enumi})\quad]
\item $\vect{AB}=\vect{E\ldots}$
\item $ \vect{AB}+\vect{FA}= \ldots\ldots$
\item $\vect{FE}+\vect{AE}= \ldots\ldots$
\item $\vect{CH}+\vect{GE} = \ldots\ldots$
\item $\vect{AE}-\vect{CD}= \ldots\ldots$
\item $\vect{AE}+\vect{CB}+\vect{EB}=\ldots\ldots$
\item $\vect{AC}+ \vect{EB}= \ldots\ldots$
\item $ \vect{BA}+ \vect{DC}-\vect{FE}= \ldots\ldots$
\end{enumerate}
%\end{multicols}
\end{minipage}

\exo

\medskip

\begin{minipage}{0.5\linewidth}
 On donne la figure ci-dessous :
 
 \medskip
 
\begin{center}
\psset{unit=1cm,algebraic=true,dimen=middle,dotstyle=o,dotsize=5pt 0,linewidth=1pt,arrowsize=3pt 2,arrowinset=0.25}
\begin{pspicture*}(-0.8,-0.8)(8.8,8.8)
\psgrid[gridlabels=0,subgriddiv=1,griddots=10](-1,-1)(9,9)
\psdots[dotstyle=*](2,1)(5,1)(1,2)(2,5)(4,3)(6,4)
\rput(1.8,0.8){$B$}
\rput(5.2,0.8){$A$}
\rput(0.8,1.8){$C$}
\rput(2.2,5.2){$D$}
\rput(3.8,2.8){$F$}
\rput(5.8,3.8){$E$}
\end{pspicture*}
\end{center}
\end{minipage}
\begin{minipage}{0.5\linewidth}
\begin{enumerate}
\item Tracer le représentant du vecteur $\vect{CD}$ d'origine E.
\item Tracer le représentant du vecteur $\vect{AB}$

 d'extrémité F.
\item Construire le vecteur $\vect{u}=\vect{AB}+\vect{CD}$
\item construire le vecteur $\vect{v}=\vect{AB}-\vect{CD}$
\end{enumerate}
\end{minipage}


\exo 

\medskip


R, S et T sont trois points non alignés.

\begin{enumerate}
\item \begin{enumerate}
	\item Construire le point P tel que $\vect{RP}=\vect{RS}+\vect{RT}$
	\item En utilisant la relation de Chasles, montrer que $\vect{TP}=\vect{RS}$
\end{enumerate}
\item \begin{enumerate}
\item Construire le point U tel que $\vect{SU}=\vect{SR} - \vect{TS}$
\item Montrer que $\vect{RU}=\vect{ST}$
\item En déduire que T est le milieu de [UP].
\end{enumerate}
\end{enumerate}



\exo

\begin{minipage}{0.5\linewidth}

 On donne la figure ci-dessous :
 
 \medskip
 
\begin{center}
\psset{unit=1cm,algebraic=true,dimen=middle,dotstyle=o,dotsize=5pt 0,linewidth=1pt,arrowsize=3pt 2,arrowinset=0.25}
\begin{pspicture*}(0.2,0.2)(8.8,8.8)
\psgrid[gridlabels=0,subgriddiv=1,griddots=10](-1,-1)(9,9)
\psdots[dotstyle=*](1,8)(3,7)(7,5)(3,1)(4,2)(7,5)
\rput(0.8,8.2){$C$}
\rput(3.2,7.2){$B$}
\rput(6.7,5.5){$A$}
\rput(2.7,1.5){$D$}
\rput(3.7,2.5){$E$}
\psline(0,8.5)(9,4)
\psline(2,0)(9,7)
\end{pspicture*}
\end{center}
\end{minipage}
\begin{minipage}{0.5\linewidth}
\begin{enumerate}
\item Exprimer le vecteur $\vect{AC}$ en fonction de $\vect{AB}$

\medskip

\dotfill

\medskip

\dotfill

\item Exprimer le vecteur $\vect{AE}$ en fonction de $\vect{DE}$

\medskip

\dotfill

\medskip

\dotfill
\end{enumerate}
\end{minipage}


\exo

\medskip

\begin{minipage}{0.5\linewidth}
 On donne la figure ci-dessous :
 
 \medskip
 
\begin{center}
\psset{unit=0.5cm,algebraic=true,dimen=middle,dotstyle=o,dotsize=5pt 0,linewidth=1pt,arrowsize=3pt 2,arrowinset=0.25}
\begin{pspicture*}(-0.5,-0.5)(15.5,9.5)
\psgrid[gridlabels=0,subgriddiv=1,griddots=5](-1,-1)(16,10)
\psline[linewidth=1.2pt]{->}(6,5)(5,1)
\psline[linewidth=1.2pt]{->}(11,6)(14,8)

\psdots[dotstyle=*](3,5)(9,2)(9,8)

\rput(12.5,7.5){$\vect{u}$}
\rput(6.2,3.5){$\vect{v}$}

\rput[br](2.8,5.2){$B$}
\rput[tl](9.2,2.2){$A$}
\rput[tl](9.2,8.2){$C$}
\end{pspicture*}
\end{center}
 
 \end{minipage}
\begin{minipage}{0.5\linewidth}
\begin{enumerate}
\item Construire le point $D$ tel que

 $\vect{AD} = 2\vect{u}+\vect{v}$
\item Construire le point $F$ tel que

 $\vect{BF}=-\vect{v}-\dfrac{1}{3}\vect{AB}$
\item Construire le point $E$ tel que

 $\vect{CE}=\vect{FB}-\vect{AB}-\vect{v}-3\vect{u}$


\end{enumerate}
%\end{multicols}
\end{minipage}

\exo

On considère un parallélogramme KLMN de centre I.

Faire une figure puis :

\begin{enumerate}
\item Construire le point H tel que $\vect{LH}=\dfrac{4}{3}\vect{KL}$
\item Construire le point D tel que $\vect{ID}=-\dfrac{3}{2}\vect{KN}$
\item Construire le point V tel que $\vect{VL}=\dfrac{5}{4}\vect{MN}$

\end{enumerate}

\exo 

ABC est un triangle quelconque.

On définit les points $D$ et $E$ par les relations vectorielles : $\vect{AD}=\dfrac{5}{3} \vect{AB}$ et $\vect{AE}=-\dfrac{1}{2}\vect{AC}$

\begin{enumerate}
\item Faire une figure et placer les points $D$ et $E$. 
\item En utilisant judicieusement la remation de Chasles, montrer que $\vect{BD}=\dfrac{2}{3}\vect{AB}$
\item Démontrer que $\vect{CE}=-\dfrac{3}{2}\vect{AC}$
\end{enumerate}

\vspace{1cm}

\exo

ABCD est un parallélogramme, $I$ et $J$ sont les points définis par $\vect{DI}=2 \vect{IC}$ et $5 \vect{BJ} = 2 \vect{AJ}$

\begin{enumerate}
\item Exprimer le vecteur $\vect{CI}$ en fonction du vecteur $\vect{CD}$ 
\item Exprimer le vecteur $\vect{BJ}$ en fonction du vecteur $\vect{BA}$
\item Faire une figure
\item Montrer, à l'aide d'une égalité vectorielle, que DIJB est un parallélogramme.
\end{enumerate}

\vspace{1cm}

\exo

ABC est un triangle. Les points $M$ et $N$ sont définis par $\vect{BM}=2 \vect{AB}$ et $\vect{MN}= 3 \vect{BC}$

\begin{enumerate}
\item Faire une figure
\item Montrer que $\vect{AM} = 3 \vect{AB}$
\item Comparer les vecteurs $\vect{AN}$ et $\vect{AB}$.
\item Que peut-on en déduire ?
\end{enumerate}

\vspace{1cm}

\exo 

ABCD est un carré. $E$ et $F$ sont les milieux de $[AD]$ et de $[AB]$, $G$ est le symétrique de $A$ par rapport à $B$.

\begin{enumerate}
\item Faire une figure
\item Exprimer le vecteur $\vect{EF}$ en fonction des vecteurs $\vect{AB}$ et $\vect{AD}$
\item Exprimer le vecteur $\vect{CG}$ en fonction des vecteurs $\vect{AB}$ et $\vect{AD}$
\item En déduire que les droites $(EF)$ et $(CG)$ sont parallèles.
\end{enumerate}


